% SPDX-License-Identifier: CC-BY-NC-ND-4.0
% !TEX root = main.tex

\section{加法圏}

\begin{definition} \label{def: IsPreadditive}
局所小圏$\mathscr{C}$は次の条件を満たすとき\term{ぜんかほうけん}{前加法圏}{preadditive category}であるという.
\begin{itemize}
    \item 任意の$A, B \in \Ob(\mathscr{C})$に対して写像$\sigma_{(A,B)} \colon \Mor_{\mathscr{C}}(A, B) \times \Mor_{\mathscr{C}}(A, B) \to \Mor_{\mathscr{C}}(A, B)$が定まり次の条件を満たす.\
    ただし任意の$f, g \in \Mor_{\mathscr{C}}(A, B)$に対して$f + g \coloneq \sigma_{(A,B)}(f,g)$とする.
    \begin{itemize}
    \item ある$0_{A,B}^{+} \in \Mor_{\mathscr{C}}(A,B)$が存在して$(\Mor_{\mathscr{C}}(A,B), \sigma_{(A,B)}, 0_{A,B} )$はアーベル群.
    \item 任意の$A, B, C \in \Ob(\mathscr{C})$と$f,g \in \Mor_{\mathscr{C}}(A, B),\ h \in \Mor_{\mathscr{C}}(B, C)$に対して$h \circ (f + g) = h \circ f + f \circ g$成り立つ.
    \item 任意の$A, B, C \in \Ob(\mathscr{C})$と$f,g \in \Mor_{\mathscr{C}}(B, C),\ h \in \Mor_{\mathscr{C}}(A, B)$に対して$(f + g) \circ h = f \circ h + g \circ h$成り立つ.
    \end{itemize}
\end{itemize}
\end{definition}

\begin{lemma} \label{lem: additive_zero_comp}
$\mathscr{C}$を零対象$Z$を持つ前加法圏とする.\
任意の$A, B, C \in \Ob(\mathscr{C})$と$g \in \Mor_{\mathscr{C}}(B, C),\ h \in \Mor_{\mathscr{C}}(C, A)$に対して次の図式が可換になる.
% https://q.uiver.app/#q=WzAsNixbMCwwLCJCIl0sWzEsMCwiQSJdLFsxLDEsIkMiXSxbMywwLCJCIl0sWzQsMCwiQSJdLFs0LDEsIkMiXSxbMSwwLCIwX3tBLEJ9XnsrfSIsMl0sWzAsMiwiZyIsMl0sWzEsMiwiMF97QSxDfV57K30iXSxbNSwzLCIwX3tDLEJ9XnsrfSJdLFs0LDMsIjBfe0EsIEJ9XnsrfSIsMl0sWzUsNCwiaCIsMl1d
\[\begin{tikzcd}
	B & A && B & A \\
	& C &&& C
	\arrow["g"', from=1-1, to=2-2]
	\arrow["{0_{A,B}^{+}}"', from=1-2, to=1-1]
	\arrow["{0_{A,C}^{+}}", from=1-2, to=2-2]
	\arrow["{0_{A, B}^{+}}"', from=1-5, to=1-4]
	\arrow["{0_{C,B}^{+}}", from=2-5, to=1-4]
	\arrow["h"', from=2-5, to=1-5]
\end{tikzcd}\]
\end{lemma}

\begin{proof}
$g \circ 0_{A,B}^{+} = g \circ (0_{A,B}^{+} + 0_{A,B}^{+}) = g \circ 0_{A,B}^{+} + g \circ 0_{A,B}^{+}$であり両辺$g \circ 0_{A,B}^{+}$を引けば$0_{A,C}^{+} = g \circ 0_{A,B}^{+}$となる.\
同様に$0_{A,B}^{+} \circ h = 0_{C,B}^{+}$となる.
\end{proof}

\begin{proposition} \label{prop: additive_zero_eq_zero_morphism}
$\mathscr{C}$を零対象$Z$を持つ前加法圏とする.\
任意の$A,B \in \Ob(\mathscr{C})$に対して次が成り立つ.
\begin{center}
    $0_{A,B}^{+} = 0_{A,B}$
\end{center}
\end{proposition}

\begin{proof}
$\Mor_{\mathscr{C}}(Z, B)$はシングルトンなので$\iexcl_B = 0_{Z,B}^{+}$である.\
\cref{lem: additive_zero_comp}により$0_{A,B} = !_A \circ 0_{Z,B}^{+} = 0_{A,B}^{+}$となる.
\end{proof}

\begin{definition} \label{def: IsAdditive}
前加法圏$\mathscr{C}$は次の条件を満たすとき\term{かほうけん}{加法圏}{additive category}であるという.
\begin{itemize}
    \item 零対象$Z \in \Ob(\mathscr{C})$が存在.
    \item 任意の$A, B \in \Ob(\mathscr{C})$に対して$D_{A,B} \colon \mathbf{2} \to \mathscr{C}$の極限$(A \times B, \pi)$が存在.
\end{itemize}
\end{definition}
